% ============================================================================
% IJPR Special Issue manuscript draft
% Special Issue: Innovations in Production Planning: Emerging Problems and
% Modern Solution Paradigms
% Taylor & Francis Interact template with APA-style references
% ============================================================================
\documentclass[]{interact}

\usepackage{epstopdf}
\usepackage[caption=false]{subfig}
\usepackage[longnamesfirst,sort]{natbib}
\bibpunct[, ]{(}{)}{;}{a}{,}{,}
\renewcommand\bibfont{\fontsize{10}{12}\selectfont}
\usepackage{amsmath,amssymb}
\usepackage{booktabs}
\usepackage{array}
\usepackage{graphicx}
\usepackage{url}

\newcommand{\placeholder}[1]{\textbf{[PLACEHOLDER: #1]}}
\newcommand{\methodname}{H-NSGA-II}
\newcommand{\Kset}{\mathcal{K}}
\newcommand{\Lset}{\mathcal{L}}
\newcommand{\Tset}{\mathcal{T}}
\newcommand{\Gset}{\mathcal{G}}
\newcommand{\Fset}{\mathcal{F}}

\begin{document}

\articletype{RESEARCH ARTICLE}

\title{Agricultural Production Decisions under Price Risk: Exact Multiobjective Optimization and Surrogate-Assisted NSGA-II}

\author{
\name{Leonardo H. Talero-Sarmiento\textsuperscript{a}\thanks{CONTACT Leonardo H. Talero-Sarmiento. Email: \placeholder{corresponding author email}}}
\affil{\textsuperscript{a}\placeholder{institution, department, city, country; add coauthors and affiliations before submission}}
}

\maketitle

\begin{abstract}
Agricultural production decisions combine crop choice, planting time, land allocation, and farm participation, while biological lead times and volatile output prices expose producers to trade-offs that cannot be represented adequately through a single economic criterion.  This study formulates a data-driven bi-objective planning problem that maximizes expected agricultural income while minimizing covariance-based price risk across heterogeneous farms, products, planting periods, and municipality-specific productivity conditions.  Three solution paradigms are evaluated through a common decision model, because an exact $\varepsilon$-constraint procedure implemented with SCIP provides reference Pareto solutions for tractable instances, NSGA-II provides a scalable evolutionary benchmark, and a surrogate-assisted hybrid NSGA-II combines evolutionary schedule search with conditional convex quadratic refinement of cultivated areas.  The hybrid method uses a learned gate to allocate mathematical-programming calls selectively, while every retained solution is evaluated with the true objectives and feasibility rules rather than surrogate predictions.  The empirical design uses 17,059 productivity observations across 17 municipalities and 40 agricultural products, together with a reproducible benchmark ladder ranging from a two-lot smallholder problem to a multi-farm network containing 150 heterogeneous lots.  \placeholder{Insert final exact, NSGA-II, H-NSGA-II, price-forecast, demand, and statistical results after the reproducible experiment is completed.}
\end{abstract}

\begin{keywords}
agricultural decision support; multiobjective optimization; NSGA-II; surrogate-assisted evolutionary optimization; matheuristics; price risk; production planning
\end{keywords}

\section{Introduction}

Agricultural production planning requires commitments on land, crop portfolios, and planting calendars before market outcomes become fully observable, while crop maturation and land-recovery requirements restrict the ability to reverse decisions after production has started.  These conditions connect agronomic feasibility with economic exposure, because local productivity determines the physical output generated by each land allocation while market prices determine the value and risk of that output across the planning horizon \citep{Ahumada2009,Tan2012,Burbano2022}.

Multiobjective crop-planning research increasingly represents decisions through mathematical programming, evolutionary optimization, or interactive decision-support procedures, particularly when economic, environmental, social, and resource objectives create structurally conflicting planning outcomes.  Recent Colombian studies have used multiobjective mathematical programming to coordinate fruit-crop decisions for small producers and to support yam cultivation and distribution planning, while other agricultural studies have extended NSGA-II with expert rules or interactive preference information to improve feasibility, computational efficiency, or decision acceptance \citep{Vergara2025,Huang2025,Arismendi2026,Kropp2026}.  This literature demonstrates that agricultural planning benefits from both prescriptive optimization and decision-support interpretation, although solution quality and computational scalability remain closely connected to the representation chosen for the underlying planning problem.

Production-planning research outside agriculture provides a complementary methodological direction, because recent matheuristics embed exact mathematical subproblems inside evolutionary or local-search procedures to exploit problem structure without requiring a complete mixed-integer solve at every search step.  Multiobjective scheduling research has combined evolutionary exploration with mathematical decoding and learning-based intensification, while integrated production and distribution studies have shown that exact and heuristic components can be coordinated effectively under strict computational budgets \citep{Hu2024,Vie2025,Kucukkoc2024}.  These developments motivate a solution architecture in which exact optimization is retained where it provides valuable local structure, while evolutionary search explores the larger combinatorial decision space.

Surrogate-assisted evolutionary optimization adds a learning layer to this architecture by predicting which candidates deserve expensive evaluation, refinement, or high-fidelity computation, thereby allocating limited computational resources toward decisions with greater expected value.  Recent reviews distinguish global, local, hybrid, and learning-based surrogate-management strategies for expensive combinatorial optimization, while newer algorithms use comparison models, adaptive local regions, and selective high-fidelity evaluation to improve convergence under restricted evaluation budgets \citep{Liu2024,Pierce2025,Wang2026}.  The relevant opportunity for agricultural planning is therefore not to replace the true objectives with machine-learning predictions, but to learn when an exact local refinement is sufficiently promising to justify its computational cost.

This study addresses that opportunity through a common bi-objective production-decision model that connects exact optimization, NSGA-II, and surrogate-assisted hybrid search under the same productivity, price, demand, and agronomic inputs.  The mathematical formulation represents price risk directly through a positive-semidefinite covariance matrix and omits product-pair auxiliary variables, while temporal feasibility is represented through precomputed occupancy and botanical-family conflict structures that avoid generic large constants.  The resulting model supports exact reference computation on smaller instances and preserves the same decision semantics when NSGA-II and the proposed \methodname{} are applied to larger farm networks.

Three research questions structure the computational experiment, with RQ1 examining how exact multiobjective tractability changes as farms, lots, products, and planning periods increase, RQ2 examining how closely standard NSGA-II reproduces exact or best-known trade-offs under comparable computational budgets, and RQ3 examining whether surrogate-guided quadratic refinement improves Pareto quality per unit of computing time.  These questions keep the analysis centered on decision accessibility, because the practical contribution depends on whether producers or coordinating organizations can obtain diverse and reliable income-risk alternatives within computational limits that remain compatible with planning practice.

\section{Related decision and solution research}

\subsection{Agricultural production decisions under multiple criteria}

Whole-farm planning under risk can be interpreted as a portfolio decision because crops generate heterogeneous returns and correlated sources of uncertainty, which makes diversification and resource allocation central to the economic interpretation of agricultural plans \citep{Radulescu2012,Burbano2022}.  Recent crop-planning applications have broadened this view by incorporating environmental, social, water, loss, and preference criteria, while mathematical programming and evolutionary algorithms continue to provide complementary ways of generating alternatives for decision-makers who face nontrivial trade-offs \citep{Vergara2025,Huang2025,Arismendi2026}.

The decision-support dimension has become increasingly explicit because computationally efficient plans are not sufficient when decision-makers cannot interpret or navigate the resulting trade-offs, particularly in settings where farm heterogeneity and local production conditions affect feasibility.  Interactive evolutionary optimization has recently been applied to agricultural irrigation decisions with fewer function evaluations than standard NSGA-II while maintaining comparable performance, which reinforces the value of linking algorithmic efficiency with decision accessibility rather than treating optimization as an isolated numerical exercise \citep{Kropp2026}.

\subsection{Matheuristics and surrogate-assisted multiobjective search}

Matheuristics combine mathematical programming with heuristic or evolutionary search by assigning each component the decisions that it can handle most effectively, which can preserve exact local structure while avoiding a monolithic combinatorial solve on large instances.  A recent multiobjective assembly-scheduling method uses mathematical decoding to optimize continuous or assignment decisions within an evolutionary procedure, while production-distribution research similarly combines exact and local-search mechanisms under realistic time restrictions \citep{Hu2024,Vie2025}.  The agricultural planning structure considered here is well suited to this logic because discrete crop schedules define a combinatorial backbone, whereas cultivated areas remain continuous decisions that can be optimized conditionally once the schedule is fixed.

Surrogate assistance becomes useful when even these conditional exact refinements are too expensive to apply indiscriminately, especially as population size and problem dimension increase across repeated generations.  Surveys and recent high-dimensional studies show that surrogate management must balance exploitation of candidates predicted to be promising with continued exploration of uncertain regions, while comparison-based and adaptive-region approaches demonstrate that learning can guide expensive multiobjective evaluation without replacing the final high-fidelity objectives \citep{Liu2024,Pierce2025,Wang2026}.  This study follows that principle by using a surrogate only to prioritize conditional QP refinement, while all nondominated sorting, archive membership, and reported decision outcomes use verified objective values.

\section{Decision setting and data architecture}

\subsection{Planning decisions}

The planning system contains a set of farms $\Fset$ partitioned into heterogeneous lots $\Lset$, where each lot may remain unused or receive one crop at a time and every accepted production decision specifies a crop, a feasible sowing period, a harvest period, and a cultivated area that cannot exceed lot capacity.  Crop-specific maturation times, botanical-family rotation rules, homogeneous preparation periods, expected prices, market demand, and municipality-dependent productivity jointly determine whether each candidate plan is feasible and economically relevant.

The decision maker evaluates each plan through expected agricultural income and portfolio price risk, while production from each crop is valued at the expected price associated with its harvest period and the covariance structure captures simultaneous exposure to correlated price movements.  Yield can vary across municipalities and planting months because the productivity layer uses empirical local evidence rather than assuming homogeneous crop performance across every farm-lot combination.

\subsection{Productivity evidence and eligibility}

The productivity dataset contains 17,059 observations across 17 municipalities and 40 agricultural products, with each record identifying municipality, product, planting month, and observed yield in kilograms per square metre.  Coverage is unbalanced across municipality-product combinations, so the computational pipeline uses hierarchical shrinkage across product, municipality-product, product-month, and municipality-product-month levels while preserving explicit support counts for every estimate used in the benchmark.

Crop eligibility is determined from empirical municipality-product support before each instance is created, which prevents the model from treating every product as equally plausible in every location when the available evidence does not support that assumption.  \placeholder{Report final support thresholds, sensitivity analysis, and the treatment of anomalous crop distributions, including the Frijol yield pattern identified during data auditing.}

\subsection{Prices, demand, and uncertainty}

Expected prices will be constructed from DANE-SIPSA wholesale market histories at weekly frequency, with the final forecasting experiment comparing ARIMA or SARIMA specifications, Gaussian-process models, and ARIMA-plus-Gaussian-process residual models through rolling-origin validation.  Conditional volatility will be evaluated separately through GARCH-family specifications when diagnostic evidence supports time-varying variance, because mean forecasting and volatility forecasting serve different roles in the production-decision model and should not be conflated within one predictive comparison.

Market demand will be obtained from SIPSA-A supply information and harmonized with the sales groups used by the production model, which provides a data-based upper bound on harvest quantities rather than an unrestricted revenue opportunity.  \placeholder{Insert final data window, market location, product mapping, missing-data treatment, rolling-origin horizons, forecast metrics, volatility model selection, and demand aggregation rules.}

\begin{table}
\tbl{Data elements required by the production-decision model and their intended empirical sources.}
{\begin{tabular}{>{\raggedright\arraybackslash}p{3.3cm}>{\raggedright\arraybackslash}p{5.0cm}>{\raggedright\arraybackslash}p{5.0cm}}
\toprule
Decision input & Empirical construction & Final evidence to report \\
\midrule
Yield & Municipality-product-month shrinkage estimates & Support, dispersion, sensitivity \\
Crop eligibility & Municipality-product empirical support & Threshold and robustness \\
Expected price & Rolling-origin selected forecast by product & MAE, RMSE, sMAPE, MASE \\
Price risk & Regularized covariance with optional conditional volatility & PSD diagnostics, stability \\
Demand & SIPSA-A weekly or harmonized group demand & Coverage and aggregation \\
Planting feasibility & Historical planting months and agronomic maturation & Calendar sensitivity \\
\bottomrule
\end{tabular}}
\label{tab:data}
\end{table}

\section{Bi-objective production-decision model}

Let $k\in\Kset$ index products, $l\in\Lset$ index lots, and $t\in\Tset$ index weekly periods, while $A_l$ denotes lot area, $r_{lks}$ denotes expected yield when product $k$ is planted in lot $l$ during period $s$, and $p_{kt}$ denotes the expected sale price when product $k$ is harvested in period $t$.  The binary decision $y_{lkt}$ equals one when product $k$ is scheduled in lot $l$ for harvest in period $t$, while $z_{lkt}\geq0$ denotes cultivated area and $x_{kt}$ denotes aggregate harvested production across all lots.

Only temporally feasible variables are created, because $\Tset_k\subseteq\Tset$ contains harvest periods that correspond to an admissible planting period and allow maturation before the planning horizon ends.  Aggregate production is defined through
\begin{equation}
 x_{kt}=\sum_{l\in\Lset} r_{lk,s(k,t)}z_{lkt},
 \qquad \forall k\in\Kset,\;t\in\Tset_k,
\label{eq:production}
\end{equation}
where $s(k,t)$ identifies the planting period associated with harvest period $t$ for product $k$ under its maturation duration.

Expected income is measured in monetary units through
\begin{equation}
 \max F_1=\sum_{k\in\Kset}\sum_{t\in\Tset_k}p_{kt}x_{kt},
\label{eq:income}
\end{equation}
while monetary exposure is defined as $e_{kt}=p_{kt}x_{kt}$ and portfolio price risk is represented directly through
\begin{equation}
 \min F_2=\sum_{t\in\Tset}\mathbf e_t^{\top}\Sigma\mathbf e_t,
\label{eq:risk}
\end{equation}
where $\Sigma$ is a symmetric positive-semidefinite covariance matrix estimated from product-price returns or from a conditional volatility and correlation decomposition selected during the forecasting stage.

Cultivated area is linked to crop activation through the physical lot bound
\begin{equation}
 0\leq z_{lkt}\leq A_l y_{lkt},
 \qquad \forall l\in\Lset,\;k\in\Kset,\;t\in\Tset_k,
\label{eq:link}
\end{equation}
which avoids an arbitrary large constant and gives the mixed-integer relaxation a bound that follows directly from the decision system.

Demand is imposed at the corresponding sales-group level when product-level demand is unavailable, with $g(k)$ denoting the sales group associated with product $k$ and $d_{gt}$ denoting the available market quantity during period $t$,
\begin{equation}
 \sum_{k:g(k)=g}x_{kt}\leq d_{gt},
 \qquad \forall g\in\Gset,\;t\in\Tset.
\label{eq:demand}
\end{equation}

Temporal occupation and rotation feasibility are represented through precomputed conflict cliques, where one set prevents simultaneous lot occupation and another expands the occupied interval for crops belonging to the same botanical family when additional recovery time is required.  This representation yields constraints of the form
\begin{equation}
 \sum_{(k,t)\in\mathcal C_{lw}}y_{lkt}\leq1,
 \qquad \forall l\in\Lset,\;w\in\Tset,
\label{eq:occupancy}
\end{equation}
and
\begin{equation}
 \sum_{(k,t)\in\mathcal C^{f}_{lw}}y_{lkt}\leq1,
 \qquad \forall l\in\Lset,\;f,\;w\in\Tset,
\label{eq:family}
\end{equation}
where the second clique family is generated only among crops sharing botanical family $f$ and therefore adds recovery time without imposing unnecessary restrictions across different families.

The resulting variable order is dominated by $O(KLT)$ crop-timing and area decisions, while the covariance objective remains quadratic rather than introducing one auxiliary variable for every product pair, lot, and period.  This compact structure is essential for the experiment because the exact method must provide credible reference fronts on at least the smaller instances before evolutionary and hybrid scalability can be interpreted meaningfully.

\begin{table}
\tbl{Structural choices used to keep the mathematical decision model compact and interpretable.}
{\begin{tabular}{>{\raggedright\arraybackslash}p{4.0cm}>{\raggedright\arraybackslash}p{4.5cm}>{\raggedright\arraybackslash}p{4.3cm}}
\toprule
Model component & Representation & Computational implication \\
\midrule
Price risk & Direct quadratic covariance term & Avoids $K^2LT$ auxiliary variables \\
Activation-area link & $z_{lkt}\leq A_l y_{lkt}$ & Removes arbitrary Big-$M$ \\
Planting calendar & Variables created only for feasible periods & Structural presolve before solver call \\
Lot occupation & Weekly conflict cliques & Stronger than disjunctive timing logic \\
Same-family recovery & Botanical-family expanded cliques & Adds only relevant rotation conflicts \\
Productivity & Municipality and planting-time dependent & Preserves heterogeneous production evidence \\
\bottomrule
\end{tabular}}
\label{tab:compact}
\end{table}

\section{Solution methods}

\subsection{Exact multiobjective optimization with SCIP}

The exact reference method uses an $\varepsilon$-constraint procedure implemented through SCIP and PySCIPOpt, because the approach can recover efficient mixed-integer solutions without relying on a single weighted aggregation of the two objectives.  A first solve maximizes $F_1$ and establishes $F_1^{\max}$, after which each frontier point minimizes the quadratic risk objective subject to $F_1\geq\varepsilon_jF_1^{\max}$ for a predefined set of aspiration fractions $\varepsilon_j\in[0,1]$.

SCIP receives the risk criterion through an epigraph variable constrained above the convex quadratic expression, while binary crop-timing decisions, continuous areas, demand limits, and conflict cliques remain in the same mathematical model.  Each $\varepsilon$ point is solved in an independent model instance, and an optional lexicographic tie-break maximizes income at the attained risk level so that the retained representative is strongly efficient whenever the requested tolerance and time limit permit certification.

The exact experiment will record solve status, relative optimality gap, wall-clock time, search nodes, and the number of efficient points recovered within the common resource limit.  \placeholder{Insert SCIP and PySCIPOpt versions, hardware, thread count, time limit per epsilon point, optimality tolerance, and the largest instance with a certified or bounded reference frontier.}

\subsection{NSGA-II benchmark}

NSGA-II is retained as the evolutionary benchmark because the decision problem contains two conflicting objectives and therefore does not require reference-direction machinery designed specifically for many-objective search \citep{Deb2002}.  Each chromosome represents crop identity, waiting time before the next feasible planting event, and cultivated-area proportion across a bounded sequence of slots for every lot, while a feasibility-aware decoder enforces eligibility, maturation, planting calendars, demand limits, setup time, and same-family recovery before objective evaluation.

The benchmark will use matched random seeds, documented parameter tuning, and equal wall-clock comparisons against the hybrid method, because comparing only generations would ignore the mathematical-programming calls executed during hybrid refinement.  \placeholder{Insert final population size, generations, crossover and mutation rates, seed count, tuning procedure, and stopping rule after the experiment is frozen.}

\subsection{Surrogate-assisted hybrid NSGA-II}

The proposed \methodname{} treats crop scheduling and cultivated-area allocation as related but computationally distinct decision layers, because NSGA-II explores the discrete schedule while a conditional convex quadratic program optimizes areas for selected schedules.  Given a decoded schedule, the refinement model minimizes price risk subject to preserving at least the candidate's pre-refinement income, lot-area bounds, and market-demand constraints, which guarantees that accepted refinements do not intentionally exchange income for lower risk inside the local mathematical search.

Applying this conditional QP to every offspring defines a full-refinement matheuristic that isolates the value created by mathematical intensification, while the surrogate-assisted method adds a learning gate that determines which offspring receive the QP call.  The implemented gate uses an ExtraTrees regression model trained on realized refinement gains, with features describing schedule size, crop diversity, botanical-family diversity, land utilization, local yields, expected prices, price volatility, farm coverage, municipality coverage, pre-refinement objectives, Pareto rank, crowding distance, and generation index.

Warm-up generations allocate refinement calls without surrogate control so that the training sample contains observed improvement outcomes, after which most of the refinement budget is assigned to candidates with the highest predicted gain while a smaller random share preserves exploration outside the learned ranking.  Surrogate predictions never replace the objective functions, because every accepted refinement and every final nondominated solution is evaluated using the mathematical production, income, risk, demand, and agronomic relationships defined by the decision model.

The computational study therefore contains four experimentally distinct configurations, namely exact $\varepsilon$-constraint optimization, standard NSGA-II, NSGA-II with full conditional QP refinement, and surrogate-gated \methodname{}, which allows the contribution of mathematical refinement to be separated from the contribution of learning-guided call allocation.  This ablation is necessary because an apparent hybrid advantage would remain ambiguous if the experiment compared only standard NSGA-II against a method that simultaneously added exact intensification and selective machine-learning guidance.

\begin{figure}
\centering
\fbox{\parbox[c][7.0cm][c]{0.88\textwidth}{\centering
\textbf{FIGURE PLACEHOLDER 1}\\[8pt]
Recommended content: end-to-end decision architecture linking productivity, SIPSA price and demand updating, forecast selection, covariance estimation, compact MIQCP, exact $\varepsilon$-constraint optimization, NSGA-II, conditional QP refinement, and surrogate gating.\\[10pt]
\placeholder{Replace with vector figure before submission.}}}
\caption{Data, decision, optimization, evolutionary-search, and surrogate-guidance architecture used in the production-planning experiment.}
\label{fig:framework}
\end{figure}

\section{Experimental design}

\subsection{Benchmark ladder}

The benchmark increases planning scale through eight reproducible instances that vary the number of farms, lots, products, aggregate area, and planning horizon, while every individual farm remains within the smallholder area context used for the empirical application.  The smaller instances are intended for exact-front benchmarking, the middle instances identify the transition where exact computation becomes increasingly expensive, and the network instances test whether evolutionary and hybrid search preserve diverse income-risk alternatives after full exact frontier construction becomes computationally restrictive.

\begin{table}
\tbl{Benchmark ladder from a small farm to a heterogeneous multi-farm production network.}
{\begin{tabular}{lrrrrl}
\toprule
Instance & Farms & Lots & Area (ha) & Products & Horizon \\
\midrule
I01 Micro & 1 & 2 & 0.6 & 5 & 52 weeks \\
I02 Small & 1 & 4 & 1.8 & 8 & 52 weeks \\
I03 Diverse & 1 & 8 & 4.5 & 12 & 104 weeks \\
I04 Two farms & 2 & 8 & 6.0 & 8 & 52 weeks \\
I05 Coop-small & 5 & 20 & 15.0 & 12 & 104 weeks \\
I06 Coop-medium & 10 & 40 & 30.0 & 12 & 104 weeks \\
I07 Network & 20 & 80 & 60.0 & 21 & 104 weeks \\
I08 Network-large & 30 & 150 & 100.0 & 21 & 104 weeks \\
\bottomrule
\end{tabular}}
\label{tab:instances}
\end{table}

Product sets are nested across the benchmark so that increasing $K$ adds decision variety without changing the complete experimental definition, while municipality assignments and heterogeneous lot areas are generated from fixed seeds and stored with each instance.  \placeholder{Confirm the final product subsets after productivity-support auditing and attach the complete instance files as supplementary material.}

\subsection{Performance measures and statistical comparison}

Exact, evolutionary, and hybrid methods will be compared under common wall-clock limits and explicitly reported evaluation counts, because the hybrid method consumes additional solver calls whose cost must be visible when decision quality is interpreted.  For instances with an exact reference set, Pareto quality will be measured through hypervolume, IGD$^+$, additive $\epsilon$-indicator, runtime, and coverage of the certified or bounded reference region, while larger instances will use a verified empirical best-known nondominated set assembled from all methods and replications.

The hybrid methods will additionally report conditional QP calls, useful-refinement hit rate, cumulative QP time, surrogate prediction error, computational savings relative to full refinement, and the share of final nondominated solutions that received mathematical intensification.  \placeholder{Use 30 independent seeds when computationally feasible, with a minimum predeclared fallback of 20 seeds if the complete SCIP-assisted experiment cannot finish before the submission cutoff, and report paired nonparametric comparisons with effect sizes.}

\begin{table}
\tbl{Evidence required to distinguish the exact, evolutionary, full-refinement, and surrogate-gated methods.}
{\begin{tabular}{>{\raggedright\arraybackslash}p{4.2cm}cccc}
\toprule
Evidence & Exact & NSGA-II & NSGA-II+QP & \methodname{} \\
\midrule
Optimality status or gap & Yes & -- & -- & -- \\
HV, IGD$^+$, $\epsilon^+$ & Reference & Yes & Yes & Yes \\
Wall-clock time & Yes & Yes & Yes & Yes \\
Repeated random seeds & -- & Yes & Yes & Yes \\
Conditional QP calls & -- & -- & All offspring & Selective \\
Refinement hit rate & -- & -- & Yes & Yes \\
Surrogate prediction evidence & -- & -- & -- & Yes \\
\bottomrule
\end{tabular}}
\label{tab:evidence}
\end{table}

\section{Results}

\subsection{Forecasting and market inputs}

\placeholder{Table 4 should report rolling-origin forecasting results by product and model, including MAE, RMSE, sMAPE, MASE, selected mean model, selected volatility model, and the forecast horizon used by each benchmark instance.}  \placeholder{Figure 2 should show representative observed and forecast price trajectories with uncertainty bands for products selected from different volatility and seasonality profiles.}

\subsection{Exact reference fronts and scalability}

\placeholder{Table 5 should report SCIP results for I01--I04, including epsilon level, income, risk, status, relative gap, nodes, and solve time, while the text should distinguish mathematically certified points from time-limited incumbents with nonzero gap.}  \placeholder{Figure 3 should overlay the exact frontier and approximate archives for at least one small and one transition-scale instance using common objective units.}

\subsection{NSGA-II and hybrid decision quality}

\placeholder{Table 6 should report seed-level summaries of HV, IGD$^+$, additive epsilon, runtime, and archive size for NSGA-II, NSGA-II+QP, and H-NSGA-II across the complete benchmark.}  \placeholder{The principal inferential comparison should identify where conditional mathematical refinement materially changes decision quality and whether surrogate gating preserves that gain with fewer solver calls.}

\subsection{Surrogate-gating ablation}

\placeholder{Table 7 should report QP calls, useful-refinement hit rate, cumulative QP time, surrogate prediction error, and final HV for full refinement versus surrogate-gated refinement.}  \placeholder{Figure 4 should plot hypervolume against cumulative QP calls and wall-clock time so that the learning contribution is visible as computational allocation rather than as an unqualified algorithmic improvement.}

\begin{table}
\tbl{Reserved structure for the principal exact, evolutionary, and hybrid comparison.}
{\begin{tabular}{llrrrrrr}
\toprule
Instance & Method & HV & IGD$^+$ & $\epsilon^+$ & Time (s) & QP calls & Gap \\
\midrule
I01 & Exact & \placeholder{} & -- & -- & \placeholder{} & -- & \placeholder{} \\
I01 & NSGA-II & \placeholder{} & \placeholder{} & \placeholder{} & \placeholder{} & -- & -- \\
I01 & NSGA-II+QP & \placeholder{} & \placeholder{} & \placeholder{} & \placeholder{} & \placeholder{} & -- \\
I01 & \methodname{} & \placeholder{} & \placeholder{} & \placeholder{} & \placeholder{} & \placeholder{} & -- \\
\multicolumn{8}{c}{\placeholder{repeat for I02--I08, with complete seed-level results in supplementary files}} \\
\bottomrule
\end{tabular}}
\label{tab:results}
\end{table}

\begin{figure}
\centering
\fbox{\parbox[c][6.5cm][c]{0.88\textwidth}{\centering
\textbf{FIGURE PLACEHOLDER 2}\\[8pt]
Recommended content: exact and approximate Pareto fronts for I01 and the largest instance where SCIP still provides a useful bounded reference, together with one large network instance represented by the empirical best-known set.\\[10pt]
\placeholder{Generate from final experiment.}}}
\caption{Income-risk decision frontiers produced by exact, evolutionary, and hybrid solution methods across representative benchmark scales.}
\label{fig:fronts}
\end{figure}

\begin{figure}
\centering
\fbox{\parbox[c][6.5cm][c]{0.88\textwidth}{\centering
\textbf{FIGURE PLACEHOLDER 3}\\[8pt]
Recommended content: computational scaling with lot-product-time size on the horizontal axis, showing SCIP time and gap together with NSGA-II and hybrid quality-per-time indicators.\\[10pt]
\placeholder{Generate from final experiment.}}}
\caption{Computational scalability and decision-quality efficiency as farms, lots, products, and planning horizon increase.}
\label{fig:scaling}
\end{figure}

\begin{figure}
\centering
\fbox{\parbox[c][6.5cm][c]{0.88\textwidth}{\centering
\textbf{FIGURE PLACEHOLDER 4}\\[8pt]
Recommended content: surrogate ablation showing hypervolume versus cumulative QP calls, refinement hit rate by generation, and full-refinement versus gated-refinement computational expenditure.\\[10pt]
\placeholder{Generate from final experiment.}}}
\caption{Effect of surrogate-guided QP allocation on mathematical-refinement cost and nondominated decision quality.}
\label{fig:surrogate}
\end{figure}

\section{Discussion}

The experiment is designed to determine which computational strategy remains decision-useful as agricultural planning scale increases, because exact optimization provides stronger optimality information while evolutionary search provides broader scalability and hybrid refinement attempts to recover mathematical structure without solving every candidate globally.  The resulting comparison should therefore be interpreted through attainable decision frontiers, computational budgets, and evidence quality rather than through a single algorithm ranking detached from the planning context.

The mathematical representation is also part of the computational contribution because direct quadratic risk, physical activation bounds, feasible-period variable generation, and conflict cliques determine how much structure remains available to both exact and hybrid solvers.  A strong formulation can postpone the point at which evolutionary approximation becomes necessary, while a conditional QP can exploit the same structure locally after a schedule has been identified by population-based search.

The surrogate layer will be considered useful only when it reduces refinement expenditure without systematically removing candidates that later contribute to the nondominated archive, which makes the full-refinement ablation and exploration safeguard essential to the interpretation.  If surrogate gating produces comparable Pareto quality with materially fewer QP calls, the learning component can be interpreted as computational resource allocation; however, if equivalent quality requires nearly complete refinement, the mathematical intensification rather than the surrogate model would constitute the meaningful hybrid contribution.

Empirical limitations must remain visible because productivity coverage is uneven across municipalities and crops, market mappings may differ across historical sources, and forecast uncertainty can influence both expected revenue and the covariance structure used by the risk objective.  Sensitivity analysis should therefore evaluate support thresholds, covariance regularization, demand aggregation, planting-calendar expansion, and forecast-model selection before managerial conclusions are drawn from particular crop or farm allocations.

\section{Conclusion}

This study formulates agricultural production planning as a multiobjective decision problem in which heterogeneous productivity, crop timing, land allocation, market demand, expected prices, and correlated price risk are evaluated through a common mathematical structure.  The computational design compares exact $\varepsilon$-constraint optimization with NSGA-II and a surrogate-assisted hybrid NSGA-II that selectively embeds convex quadratic refinement, thereby linking certified reference computation, scalable evolutionary exploration, and learning-guided mathematical intensification within one decision-support experiment.

\placeholder{Replace this paragraph with final evidence on exact tractability, NSGA-II approximation quality, QP-refinement value, surrogate savings, and the practical scale at which each method provides the most useful production decisions.}  The final interpretation should emphasize how method choice changes the accessibility and reliability of income-risk alternatives for single farms, cooperatives, and larger coordinated production networks.

\section*{Disclosure statement}

\placeholder{Add the journal-required conflict-of-interest and generative-AI disclosure statements after confirming the final author declarations and Taylor \& Francis policy wording.}

\section*{Data availability statement}

The reproducibility package will contain productivity preprocessing, price and demand updating, forecast evaluation, benchmark generation, SCIP exact optimization, NSGA-II, conditional QP refinement, surrogate-assisted \methodname{}, validation scripts, and all data permitted for redistribution by their source licenses.  \placeholder{Insert repository URL and DOI after the final computational snapshot is archived.}

\section*{Funding}

\placeholder{Add funding information or state that the research received no specific grant from public, commercial, or not-for-profit funding agencies.}

\begin{thebibliography}{}

\bibitem[Ahumada \& Villalobos(2009)]{Ahumada2009}
Ahumada, O., \& Villalobos, J. R. (2009). Application of planning models in the agri-food supply chain: A review. \emph{European Journal of Operational Research, 196}(1), 1--20. https://doi.org/10.1016/j.ejor.2008.02.014

\bibitem[Arismendi Joaqui, Henao Villabon, \& Paredes Rodriguez(2026)]{Arismendi2026}
Arismendi Joaqui, M. C., Henao Villabon, A. J., \& Paredes Rodriguez, A. M. (2026). Multi-objective optimization of the sustainable fruit crop planning problem with a focus on small Colombian producers. \emph{Kuwait Journal of Science, 53}(3), 100596. https://doi.org/10.1016/j.kjs.2026.100596

\bibitem[Battaia \& Dolgui(2022)]{BattaiaDolgui2022}
Battaia, O., \& Dolgui, A. (2022). Hybridizations in line balancing problems: A comprehensive review on new trends and formulations. \emph{International Journal of Production Economics, 250}, 108673. https://doi.org/10.1016/j.ijpe.2022.108673

\bibitem[Bollerslev(1986)]{Bollerslev1986}
Bollerslev, T. (1986). Generalized autoregressive conditional heteroskedasticity. \emph{Journal of Econometrics, 31}(3), 307--327. https://doi.org/10.1016/0304-4076(86)90063-1

\bibitem[Deb, Pratap, Agarwal, \& Meyarivan(2002)]{Deb2002}
Deb, K., Pratap, A., Agarwal, S., \& Meyarivan, T. (2002). A fast and elitist multi-objective genetic algorithm: NSGA-II. \emph{IEEE Transactions on Evolutionary Computation, 6}(2), 182--197. https://doi.org/10.1109/4235.996017

\bibitem[Hu, Zhang, Wang, Zhang, \& Tang(2024)]{Hu2024}
Hu, Y., Zhang, L., Wang, Q., Zhang, Z., \& Tang, Q. (2024). A matheuristic-based multi-objective evolutionary algorithm for flexible assembly jobs shop scheduling problem in cellular manufacture. \emph{Swarm and Evolutionary Computation, 87}, 101549. https://doi.org/10.1016/j.swevo.2024.101549

\bibitem[Huang, Pu, Zhang, Wang, Yang, \& Gu(2025)]{Huang2025}
Huang, Z., Pu, Y., Zhang, Q., Wang, Y., Yang, J., \& Gu, Y. (2025). A Fuzzy-Expert enhanced NSGA-II approach for sustainable agricultural systems. \emph{Scientific Reports, 15}, 34070. https://doi.org/10.1038/s41598-025-14488-w

\bibitem[Kropp, Nejadhashemi, Jha, Spinner, \& Toscano Pulido(2026)]{Kropp2026}
Kropp, I., Nejadhashemi, A. P., Jha, P. K., Spinner, E., \& Toscano Pulido, G. (2026). Integrating interactive decision making into evolutionary multiobjective agricultural optimization. \emph{Computers and Electronics in Agriculture, 242}, 111359. https://doi.org/10.1016/j.compag.2025.111359

\bibitem[Kucukkoc, Aydin Keskin, Karaoglan, \& Karadag(2024)]{Kucukkoc2024}
Kucukkoc, I., Aydin Keskin, G., Karaoglan, A. D., \& Karadag, S. (2024). A hybrid discrete differential evolution--genetic algorithm approach with a new batch formation mechanism for parallel batch scheduling considering batch delivery. \emph{International Journal of Production Research, 62}(1--2), 460--482. https://doi.org/10.1080/00207543.2023.2233626

\bibitem[Burbano-Figueroa et al.(2022)]{Burbano2022}
Burbano-Figueroa, O., Sierra-Monroy, A., David-Hinestroza, A., Whitney, C., Borgemeister, C., \& Luedeling, E. (2022). Farm-planning under risk: An application of decision analysis and portfolio theory for the assessment of crop diversification strategies in horticultural systems. \emph{Agricultural Systems, 199}, 103409. https://doi.org/10.1016/j.agsy.2022.103409

\bibitem[Liu, Wang, Peng, \& Yao(2024)]{Liu2024}
Liu, S., Wang, H., Peng, W., \& Yao, W. (2024). Surrogate-assisted evolutionary algorithms for expensive combinatorial optimization: A survey. \emph{Complex \& Intelligent Systems, 10}, 5933--5949. https://doi.org/10.1007/s40747-024-01465-5

\bibitem[Pierce, Kim, \& Bazarov(2025)]{Pierce2025}
Pierce, C. M., Kim, Y.-K., \& Bazarov, I. (2025). A comparison-relationship-surrogate evolutionary algorithm for multi-objective optimization. \emph{Swarm and Evolutionary Computation, 95}, 101947. https://doi.org/10.1016/j.swevo.2025.101947

\bibitem[Radulescu, Radulescu, \& Zbaganu(2012)]{Radulescu2012}
Radulescu, M., Radulescu, C. Z., \& Zbaganu, G. (2012). A portfolio theory approach to crop planning under environmental constraints. \emph{Annals of Operations Research, 219}(1), 243--264. https://doi.org/10.1007/s10479-011-0902-7

\bibitem[Rasmussen \& Williams(2006)]{Rasmussen2006}
Rasmussen, C. E., \& Williams, C. K. I. (2006). \emph{Gaussian processes for machine learning}. MIT Press.

\bibitem[Tan \& Comden(2012)]{Tan2012}
Tan, B., \& Comden, N. (2012). Agricultural planning of annual plants under demand, maturation, harvest, and yield risk. \emph{European Journal of Operational Research, 220}(2), 539--549. https://doi.org/10.1016/j.ejor.2012.02.005

\bibitem[Vergara Rodriguez, Mendoza-Ortega, Marmol Barriosnuevo, \& Diaz Sierra(2025)]{Vergara2025}
Vergara Rodriguez, C. J., Mendoza-Ortega, G. P., Marmol Barriosnuevo, M., \& Diaz Sierra, M. (2025). Multi-objective optimization for the sustainable planning of yam cultivation and distribution in the Colombian Caribbean. \emph{Frontiers in Sustainability, 6}, 1653365. https://doi.org/10.3389/frsus.2025.1653365

\bibitem[Vie, Zufferey, \& Coelho(2025)]{Vie2025}
Vie, M.-S., Zufferey, N., \& Coelho, L. C. (2025). A production and distribution scheduling matheuristic for reducing supply chain variations. \emph{Transportation Research Part E: Logistics and Transportation Review, 194}, 103905. https://doi.org/10.1016/j.tre.2024.103905

\bibitem[Wang et al.(2026)]{Wang2026}
Wang, Q., Li, H., Zhang, W., Zhang, Y., Gong, D., Chu, F., Mohamed, A. W., \& Muhammad, I. (2026). Surrogate-assisted evolutionary algorithm with adaptive local region search for high-dimensional expensive multi-objective optimization problems. \emph{Swarm and Evolutionary Computation, 100}, 102232. https://doi.org/10.1016/j.swevo.2025.102232

\end{thebibliography}

\appendix

\section{Recommended submission figures and tables}

The main manuscript should prioritize four figures that directly support the decision and method narrative, namely the integrated framework, representative income-risk frontiers, computational scalability, and the surrogate-refinement ablation that reveals whether learning reduces mathematical-programming expenditure.  Forecast examples, productivity-support heatmaps, municipality-product coverage, complete fronts for every seed, and detailed solver logs should move to supplementary material unless one becomes necessary to justify a central modeling assumption.

The principal tables should report data coverage, benchmark instances, compact-model structure, forecasting validation, solver and algorithm settings, exact-versus-approximate performance, and the full-refinement versus surrogate-gated ablation.  Seed-level metrics, complete event schedules, instance parameter files, SCIP logs, and surrogate training diagnostics should remain in the reproducibility package so that the main text can emphasize decision interpretation rather than secondary implementation detail.

\section{Submission-critical placeholders}

Before submission, the manuscript requires confirmed authorship and affiliations, current SIPSA prices and demand, rolling-origin forecast selection, the SCIP exact benchmark, repeated NSGA-II and hybrid experiments, surrogate ablation, final figures, repository metadata, funding information, and journal-specific disclosures.  Any placeholder that remains unresolved at submission should be removed together with the corresponding unsupported claim, because a narrower fully evidenced analysis is preferable to a broader narrative that implies results which the computational experiment has not yet generated.

\end{document}
